% =====================================================================
%  CSE499A — EHR Pre-Consultation Medical Documentation System
%  Standalone section: Social and Environmental Effects
%  Authors: Rafiur Rahman Mashrafi (2221971042),
%           M.G. Rabbi Hossen (2222516042),
%           Israt Zaman Srity (2211084042)
%  Advisor: Dr. Mohammad Ashrafuzzaman Khan [AzK]
%  Department of Electrical and Computer Engineering,
%  North South University · Spring 2026
% =====================================================================

\documentclass[12pt,a4paper]{article}

\usepackage[utf8]{inputenc}
\usepackage[T1]{fontenc}
\usepackage{mathptmx}
\usepackage[a4paper,margin=1in]{geometry}
\usepackage{setspace}
\usepackage{enumitem}
\usepackage[hidelinks]{hyperref}
\usepackage{titlesec}
\usepackage{xcolor}

\onehalfspacing
\setlength{\parskip}{6pt}
\setlength{\parindent}{1.5em}

\titleformat{\section}{\normalfont\large\bfseries}{\thesection}{0.7em}{}
\titleformat{\subsection}{\normalfont\normalsize\bfseries}{\thesubsection}{0.7em}{}

\title{\bfseries Social and Environmental Effects of the\\
EHR-Based Pre-Consultation Medical Documentation System}
\author{Rafiur Rahman Mashrafi (ID 2221971042)\\
        M.~G.~Rabbi Hossen (ID 2222516042)\\
        Israt Zaman Srity (ID 2211084042)\\[6pt]
        \emph{Faculty Advisor:} Dr. Mohammad Ashrafuzzaman Khan [AzK]\\
        Department of Electrical and Computer Engineering\\
        North South University, Bangladesh}
\date{Senior Design Project (CSE499A) \quad Spring 2026}

\begin{document}
\maketitle

\section{Introduction}

This section examines the social and environmental effects of the
\textit{EHR-Based Pre-Consultation Medical Documentation System}, a
voice-driven pipeline that turns a patient's spoken description of
their symptoms in Bangla --- including regional dialects and
Bangla--English code-mixing --- into a structured Electronic Health
Record (EHR) before the consultation begins. Because the system
intersects healthcare, language, and AI, its effects extend well
beyond a single clinic: it touches how Bangladeshi patients
communicate with the formal medical system, how doctors spend their
limited consultation time, and how computational resources are used
to serve a low-resource language. The discussion below treats social
and environmental impact separately, then closes with the project's
alignment to the United Nations Sustainable Development Goals (SDGs).

\section{Social Effects}

\subsection{Reclaiming consultation time for clinical care}
The Bangladeshi public-health system operates with a chronic
imbalance between patient demand and physician supply. According to
the World Health Organization, the country has roughly seven
physicians per 10{,}000 population~\cite{rep:who}. In a typical
outpatient consultation a substantial portion of the limited time is
spent on history-taking and handwritten note preparation rather than
on examination and decision-making. By absorbing the documentation
work into the waiting period --- the patient describes their history
to a tablet or kiosk before meeting the doctor --- the proposed
system returns those minutes to actual clinical reasoning. The
projected effect is therefore not ``faster consultations,'' which
would be undesirable, but \emph{better} consultations of the same
duration: the doctor walks in already knowing the patient's story.

\subsection{Voice-first inclusion of low-literacy and rural patients}
A second, often-overlooked social effect concerns who can use a
digital health tool at all. The majority of patients outside Dhaka
cannot type fluently in Bangla on a touchscreen, and most cannot use
English-language health applications. Existing Bangladeshi
``digital health'' offerings --- when they exist --- typically assume
literacy and English fluency, and therefore exclude precisely the
populations who most need primary-care support. A voice-first design
inverts this assumption. The patient speaks in whatever Bangla they
naturally use; the system meets them halfway. Crucially, the project's
empirical commitment to dialect coverage (Puran Dhaka, Barishali,
Sylheti, Normal Bangla, Indian Bangla, with Chittagonian and
Rangpuri planned for the next phase) means that the system does not
\emph{force} patients to ``correct'' their speech to a textbook
register. This is a small but meaningful act of dignity: rural and
dialect-speaking patients are addressed by the technology in their
own language, not in a sanitized standard variety.

\subsection{Impact on physicians}
Physicians are themselves a stakeholder group whose work the system
will alter. The intended effect is reduction of clerical burden:
fewer minutes of manual note-taking, fewer handwriting errors, and a
durable structured record that can be retrieved at follow-up visits.
The system is explicitly assistive --- it never issues a diagnosis
and never auto-confirms a prescription --- so the physician's
clinical authority is unchanged. The original transcript is always
displayed alongside the extracted entities, so the physician can
verify or correct any misrecognition at a glance. This design choice
is deliberate: a system that hides what it understood from the
physician would erode trust quickly, particularly in a setting
where the doctor is professionally and legally accountable for the
final decision.

\subsection{Cultural impact and Bangla-language preservation}
On a longer horizon, the project contributes to a small but real
cultural good: it puts dialect Bangla into AI infrastructure. Most
existing Bangla speech and language models are trained on a single
``standard'' register drawn from news, audiobooks, and read sentences.
The dialects spoken by tens of millions of Bangladeshis are
under-represented in this AI infrastructure or actively normalized
out. By collecting, modeling, and reporting separately on each
regional variety, the project legitimizes those varieties as
first-class objects of computational study. The dataset and
benchmark scripts produced by the project are publicly hosted, so the
contribution outlives the immediate clinical use case.

\subsection{Risks to be managed}
Any system that touches sensitive health information carries
risks that must be actively managed rather than assumed away.
Three are particularly relevant.

\begin{enumerate}[leftmargin=*]
\item \textbf{Privacy and consent.} Audio recordings of patient
speech are highly identifying personal data. The deployment plan
requires explicit, recorded informed consent in the patient's own
language before any audio capture, plus opt-out at any time without
penalty. Audio is encrypted at rest and in transit, and the deployed
infrastructure will keep data on local servers wherever feasible.

\item \textbf{Accuracy and harm from errors.} An ASR
mis-transcription that propagates into the EHR could in principle
mislead a physician. The system mitigates this in two ways: the
original audio transcript is always shown beside the extracted
record, and the system never recommends a prescription. Reported
accuracy figures (Word Error Rate, Character Error Rate, BLEU) are
computed on real dialect speech, not on clean read-aloud benchmarks,
so the limitations are visible to the physician.

\item \textbf{Equity across dialects.} A model that performs well on
``standard Bangla'' but poorly on Sylheti or Chittagonian would
amplify rather than reduce the urban--rural digital divide.
The benchmark methodology in the present project is designed to
detect such gaps; the planned fine-tuning phase will explicitly
balance per-dialect representation in the training data.
\end{enumerate}

\section{Environmental Effects}

\subsection{Computational footprint of CSE499A}
The dominant environmental impact of an AI project is the
electricity consumed during model training and large-scale inference.
For CSE499A, all twelve baseline ASR models and the six larger
multimodal audio language models were used in \emph{inference mode
only}; no end-to-end training run was conducted, and no model was
trained from scratch. The cumulative GPU consumption across the two
benchmarks is on the order of tens of GPU-hours on free-tier Google
Colab T4 and A100 instances. This is a deliberately small footprint:
the goal of CSE499A was empirical model selection, not training.

\subsection{Choice of a mid-sized model for fine-tuning}
The empirical headline finding from CSE499A --- that 1.7B--7B
parameter multimodal audio language models underperform a smaller
Bangla-specialized baseline by more than 50 percentage points of WER
--- has a direct sustainability benefit. It justifies selecting a
mid-sized model (Qwen3-ASR-1.7B) for the planned fine-tuning rather
than a 7B-parameter audio LLM. A smaller model trains faster,
performs inference faster at deployment, and consumes substantially
less energy across its operational lifetime.

\subsection{Parameter-efficient fine-tuning}
The planned fine-tuning approach in CSE499B will use Low-Rank
Adaptation (LoRA), which freezes the base model's weights and trains
only small low-rank adapter matrices. LoRA reduces the compute and
memory footprint of fine-tuning by approximately one to two orders
of magnitude versus full fine-tuning, with proportional reduction in
electricity consumption. This is a substantive --- not cosmetic ---
sustainability decision: full fine-tuning of a 1.7B-parameter audio
model on the project's compute budget would not be feasible, but a
LoRA adapter run on a single A100 instance is.

\subsection{Edge inference at deployment}
At deployment, the trained model is intended to run on modest
clinic-side hardware: a single mid-range GPU, or a quantized CPU
build for clinics without GPU access. This avoids continuous
high-power cloud inference for every patient interaction. The
patient-facing capture device is a consumer tablet or kiosk; no
custom hardware is fabricated. From a lifecycle perspective, the
project therefore consumes no new manufacturing energy beyond what
already exists in commodity tablets.

\subsection{Paper reduction}
A small but real environmental benefit is paper reduction. The
current state of small Bangladeshi clinics is paper-heavy:
prescriptions, hand-written notes, paper patient charts. A digital
EHR pipeline displaces a portion of this paper consumption.

\section{Alignment with the UN Sustainable Development Goals}

The project supports the following SDGs.

\begin{itemize}[leftmargin=*]
\item \textbf{SDG 3 --- Good Health and Well-being.} Improving the
quality of primary-care consultations and introducing structured
health records to clinics that currently have none directly serves
this goal.
\item \textbf{SDG 4 --- Quality Education.} The dialect dataset and
the open-source code, evaluation scripts, and benchmark results
contribute to educational and research resources for Bangla speech
and clinical NLP.
\item \textbf{SDG 9 --- Industry, Innovation and Infrastructure.} The
project develops domestic AI capacity for a sector (digital health)
that is currently dependent on foreign English-only tools.
\item \textbf{SDG 10 --- Reduced Inequalities.} A voice-first,
dialect-aware system explicitly serves rural and low-literacy
populations under-served by existing digital health offerings.
\item \textbf{SDG 13 --- Climate Action.} The project's preference
for small specialized models and parameter-efficient fine-tuning is
a concrete (if modest) reduction in AI energy consumption per unit
of clinical value delivered.
\end{itemize}

\section{Conclusion}

The proposed system's social effects are strongest where Bangladesh's
existing healthcare infrastructure is weakest: in small clinics,
among low-literacy and dialect-speaking patients, and during the
short consultation windows that a typical primary-care visit
allows. Its environmental footprint, by deliberate design, is small
relative to comparable AI projects: inference-only evaluation in
CSE499A, a mid-sized rather than 7B-parameter target model for
fine-tuning, parameter-efficient adapter training, and edge
inference at deployment. Both effects align with the broader UN
sustainability agenda, particularly SDGs 3, 9, 10, and 13. The risks
that genuinely deserve management --- privacy, accuracy, dialect
equity --- are addressed through specific design choices (consent,
encryption, transcript visibility, balanced dialect data) rather
than through generic disclaimers.

% =====================================================================
% REFERENCES
% =====================================================================
\begin{thebibliography}{99}

\bibitem{rep:who}
World Health Organization,
``Density of Physicians (per 1000 population), Bangladesh,'' 2023.
\url{https://data.who.int/countries/050}.
Online; accessed April 2026.

\bibitem{rep:radford}
A.~Radford, J.~W.~Kim, T.~Xu, G.~Brockman, C.~McLeavey,
and I.~Sutskever,
``Robust Speech Recognition via Large-Scale Weak Supervision,''
in \emph{Proc.\ 40th Int.\ Conf.\ on Machine Learning (ICML)},
pp.~28492--28518, 2023.
\url{https://arxiv.org/abs/2212.04356}.

\bibitem{rep:mms}
V.~Pratap \emph{et al.},
``Scaling Speech Technology to 1{,}000+ Languages,''
\emph{arXiv preprint arXiv:2305.13516}, 2023.
\url{https://arxiv.org/abs/2305.13516}.

\bibitem{rep:lora}
E.~J.~Hu \emph{et al.},
``LoRA: Low-Rank Adaptation of Large Language Models,''
in \emph{Int.\ Conf.\ on Learning Representations (ICLR)}, 2022.
\url{https://arxiv.org/abs/2106.09685}.

\bibitem{rep:oodspeech}
F.~R.~Rakib \emph{et al.},
``OOD-Speech: A Large Bengali Speech Recognition Dataset for
Out-of-Distribution Benchmarking,''
in \emph{Proc.\ INTERSPEECH 2023}.
\url{https://arxiv.org/abs/2305.09688}.

\bibitem{rep:un}
United Nations,
``The 17 Goals --- Sustainable Development,''
United Nations Department of Economic and Social Affairs.
\url{https://sdgs.un.org/goals}.
Online; accessed April 2026.

\end{thebibliography}

\end{document}
